\documentclass[10pt,a4paper]{article}

%double si on veut une version prof et une version élève. simple si on veut une seule version.
\newcommand{\buildmode}{simple}

\InputIfFileExists{version.tex}{}{}
% Version (définie par le script)
\providecommand{\version}{eleve}

\usepackage{feuilleinterro}
\usetikzlibrary{babel}

%%%à modifier à chaque fois

\renewcommand{\niveau}{2nde}
\renewcommand{\chapitre}{Fonctions : image, antécédent, ensemble de définition}
\renewcommand{\numchap}{0.4}
\renewcommand{\theme}{Fonctions}

% Repère quadrillé + courbe (ligne brisée) de la fonction f.
% #1 : xmin, #2 : xmax, #3 : ymin, #4 : ymax (graduations entières)
% #5 : sommets de la ligne brisée, ex. {(-3,-2)(-1,2)(1,-2)}
\newcommand{\courbef}[5]{
\begin{tikzpicture}
\begin{axis}[
    axis lines=middle,
    xmin=#1-0.5, xmax=#2+0.5,
    ymin=#3-0.5, ymax=#4+0.5,
    grid=major,
    xtick={#1,...,#2},
    ytick={#3,...,#4},
    xlabel={\(x\)},
    ylabel={\(y\)},
    tick label style={font=\small},
    width=9cm,
    height=8.8cm
]
\addplot[blue, very thick] coordinates {#5};
\end{axis}
\end{tikzpicture}
}

\renewcommand{\arraystretch}{1.8}

\begin{document}

% =========================================================
% PAGE 1 : VERSION A
% =========================================================

\interroversion{A}

\textbf{Exercice 1 -- Calcul littéral \hfill /3}

Développer puis réduire l'expression suivante, en détaillant les étapes :
\[
A=3(2x-5)-4(x-2).
\]

\vspace{3.5cm}

\textbf{Exercice 2 -- Lecture graphique \hfill /5}

On considère une fonction \(f\) dont la représentation graphique est donnée ci-dessous. \textit{Laisser apparents les traits de lecture.}

\noindent
\begin{minipage}[t]{0.52\textwidth}
\vspace{0pt}
\courbef{-3}{4}{-3}{3}{(-3,-2)(-1,2)(1,-2)(3,2)(4,1)}
\end{minipage}
\hfill
\begin{minipage}[t]{0.45\textwidth}
\vspace{0.5cm}
\begin{enumerate}[leftmargin=*]
    \item Quel est l'ensemble de définition de \(f\) ?
    \vspace{1.4cm}
    \item Donner les images de \(1\) et de \(4\) par \(f\).
    \vspace{1.9cm}
    \item Donner les éventuels antécédents de \(2\) et de \(-3\) par \(f\).
    \vspace{1.9cm}
    \item Résoudre l'équation \(f(x)=0\).
\end{enumerate}
\end{minipage}

\vspace{0.8cm}

\textbf{Exercice 3 -- Tableau de valeurs \hfill /2}

On considère une fonction \(g\) définie sur l'intervalle \([-3\ ;\ 3]\) par le tableau de valeurs suivant :
\begin{center}
\begin{tabular}{|*{8}{c|}}
    \hline
    \(x\) & \(-3\) & \(-2\) & \(-1\) & \(0\) & \(1\) & \(2\) & \(3\) \\
    \hline
    \(g(x)\) & \(4\) & \(-1\) & \(2\) & \(0\) & \(-1\) & \(3\) & \(2\) \\
    \hline
\end{tabular}
\end{center}

\begin{enumerate}
    \item Quelle est l'image de \(1\) par \(g\) ?
    \vspace{1.2cm}
    \item Résoudre l'équation \(g(x)=2\).
\end{enumerate}

\newpage

% =========================================================
% PAGE 2 : VERSION B
% =========================================================

\interroversion{B}

\textbf{Exercice 1 -- Calcul littéral \hfill /3}

Développer puis réduire l'expression suivante, en détaillant les étapes :
\[
B=5(x+2)-2(3x-4).
\]

\vspace{3.5cm}

\textbf{Exercice 2 -- Lecture graphique \hfill /5}

On considère une fonction \(f\) dont la représentation graphique est donnée ci-dessous. \textit{Laisser apparents les traits de lecture.}

\noindent
\begin{minipage}[t]{0.52\textwidth}
\vspace{0pt}
\courbef{-2}{4}{-2}{3}{(-2,3)(0,-1)(2,3)(4,1)}
\end{minipage}
\hfill
\begin{minipage}[t]{0.45\textwidth}
\vspace{0.5cm}
\begin{enumerate}[leftmargin=*]
    \item Quel est l'ensemble de définition de \(f\) ?
    \vspace{1.4cm}
    \item Donner les images de \(0\) et de \(3\) par \(f\).
    \vspace{1.9cm}
    \item Donner les éventuels antécédents de \(1\) et de \(-2\) par \(f\).
    \vspace{1.9cm}
    \item Résoudre l'équation \(f(x)=3\).
\end{enumerate}
\end{minipage}

\vspace{0.8cm}

\textbf{Exercice 3 -- Tableau de valeurs \hfill /2}

On considère une fonction \(g\) définie sur l'intervalle \([-2\ ;\ 4]\) par le tableau de valeurs suivant :
\begin{center}
\begin{tabular}{|*{8}{c|}}
    \hline
    \(x\) & \(-2\) & \(-1\) & \(0\) & \(1\) & \(2\) & \(3\) & \(4\) \\
    \hline
    \(g(x)\) & \(3\) & \(0\) & \(5\) & \(\dfrac{1}{2}\) & \(0\) & \(-2\) & \(1\) \\
    \hline
\end{tabular}
\end{center}

\begin{enumerate}
    \item Déterminer \(g(1)\).
    \vspace{1.2cm}
    \item Résoudre l'équation \(g(x)=0\).
\end{enumerate}

\newpage

% =========================================================
% PAGE 3 : VERSION C
% =========================================================

\interroversion{C}

\textbf{Exercice 1 -- Calcul littéral \hfill /3}

Développer puis réduire l'expression suivante, en détaillant les étapes :
\[
C=4(3x-1)-3(2x+5).
\]

\vspace{3.5cm}

\textbf{Exercice 2 -- Lecture graphique \hfill /5}

On considère une fonction \(f\) dont la représentation graphique est donnée ci-dessous. \textit{Laisser apparents les traits de lecture.}

\noindent
\begin{minipage}[t]{0.52\textwidth}
\vspace{0pt}
\courbef{-4}{2}{-2}{3}{(-4,-1)(-2,3)(1,0)(2,1)}
\end{minipage}
\hfill
\begin{minipage}[t]{0.45\textwidth}
\vspace{0.5cm}
\begin{enumerate}[leftmargin=*]
    \item Quel est l'ensemble de définition de \(f\) ?
    \vspace{1.4cm}
    \item Donner les images de \(-4\) et de \(-1\) par \(f\).
    \vspace{1.9cm}
    \item Donner les éventuels antécédents de \(3\) et de \(-2\) par \(f\).
    \vspace{1.9cm}
    \item Résoudre l'équation \(f(x)=1\).
\end{enumerate}
\end{minipage}

\vspace{0.8cm}

\textbf{Exercice 3 -- Tableau de valeurs \hfill /2}

On considère une fonction \(g\) définie sur l'intervalle \([-1\ ;\ 5]\) par le tableau de valeurs suivant :
\begin{center}
\begin{tabular}{|*{8}{c|}}
    \hline
    \(x\) & \(-1\) & \(0\) & \(1\) & \(2\) & \(3\) & \(4\) & \(5\) \\
    \hline
    \(g(x)\) & \(2\) & \(-3\) & \(4\) & \(1\) & \(2\) & \(0\) & \(2\) \\
    \hline
\end{tabular}
\end{center}

\begin{enumerate}
    \item Compléter : \(g(4)=\ldots\ldots\)
    \vspace{1.2cm}
    \item Quels sont les nombres qui ont pour image \(2\) par \(g\) ?
\end{enumerate}

\newpage

% =========================================================
% PAGE 4 : VERSION D
% =========================================================

\interroversion{D}

\textbf{Exercice 1 -- Calcul littéral \hfill /3}

Développer puis réduire l'expression suivante, en détaillant les étapes :
\[
D=2(4x+3)-5(x-3).
\]

\vspace{3.5cm}

\textbf{Exercice 2 -- Lecture graphique \hfill /5}

On considère une fonction \(f\) dont la représentation graphique est donnée ci-dessous. \textit{Laisser apparents les traits de lecture.}

\noindent
\begin{minipage}[t]{0.52\textwidth}
\vspace{0pt}
\courbef{-1}{5}{-2}{3}{(-1,-2)(1,2)(4,-1)(5,0)}
\end{minipage}
\hfill
\begin{minipage}[t]{0.45\textwidth}
\vspace{0.5cm}
\begin{enumerate}[leftmargin=*]
    \item Quel est l'ensemble de définition de \(f\) ?
    \vspace{1.4cm}
    \item Donner les images de \(-1\) et de \(2\) par \(f\).
    \vspace{1.9cm}
    \item Donner les éventuels antécédents de \(2\) et de \(3\) par \(f\).
    \vspace{1.9cm}
    \item Résoudre l'équation \(f(x)=0\).
\end{enumerate}
\end{minipage}

\vspace{0.8cm}

\textbf{Exercice 3 -- Tableau de valeurs \hfill /2}

On considère une fonction \(g\) définie sur l'intervalle \([-2\ ;\ 4]\) par le tableau de valeurs suivant :
\begin{center}
\begin{tabular}{|*{8}{c|}}
    \hline
    \(x\) & \(-2\) & \(-1\) & \(0\) & \(1\) & \(2\) & \(3\) & \(4\) \\
    \hline
    \(g(x)\) & \(-1\) & \(\dfrac{3}{2}\) & \(4\) & \(-1\) & \(2\) & \(4\) & \(0\) \\
    \hline
\end{tabular}
\end{center}

\begin{enumerate}
    \item Quelle est l'image de \(-1\) par \(g\) ?
    \vspace{1.2cm}
    \item Le nombre \(4\) a-t-il des antécédents par \(g\) ? Si oui, lesquels ?
\end{enumerate}

\end{document}
